\documentclass[spanish,11pt,a4paper]{article}
\usepackage[margin=1in]{geometry}
\usepackage{fontspec}
\setmainfont{lmroman10-regular.otf}[BoldFont=lmroman10-bold.otf,ItalicFont=lmroman10-italic.otf,BoldItalicFont=lmroman10-bolditalic.otf]
\setsansfont{lmsans10-regular.otf}[BoldFont=lmsans10-bold.otf]
\setmonofont{lmmono10-regular.otf}[Scale=MatchLowercase]
\newlength{\LeanCharWidth}
\usepackage{amsmath,amssymb}
\usepackage[spanish,es-noshorthands,es-tabla]{babel}
\usepackage{microtype,parskip}
\usepackage{graphicx,float,caption}
\usepackage{longtable,booktabs,array,calc}
\usepackage{xcolor}
\usepackage{xurl}
\usepackage{hyperref,bookmark}
\usepackage{etoolbox,needspace}
\setcounter{secnumdepth}{-\maxdimen}
\setlength{\emergencystretch}{3em}
\providecommand{\tightlist}{\setlength{\itemsep}{0pt}\setlength{\parskip}{0pt}}
\newcounter{none}
\AtBeginEnvironment{longtable}{\Needspace{8\baselineskip}\small}
\pretocmd{\section}{\Needspace{5\baselineskip}}{}{}
\pretocmd{\subsection}{\Needspace{4\baselineskip}}{}{}
\pretocmd{\subsubsection}{\Needspace{4\baselineskip}}{}{}
\captionsetup{font=small,labelfont=bf}
\urlstyle{same}

\widowpenalty=10000
\clubpenalty=10000
\hypersetup{pdftitle={Companion formal y computacional del Paper V — Preprint 1.0},pdfauthor={Juan Pablo Traverso Gianini},pdflang={es-ES},colorlinks=true,linkcolor=blue,urlcolor=blue,citecolor=blue}
\title{Companion formal y computacional del Paper V\\Preprint 1.0}
\author{Juan Pablo Traverso Gianini}
\date{18 de octubre de 2026}
\begin{document}
\raggedbottom
\maketitle
\setlength{\parskip}{5pt}
Este documento describe el desarrollo formal final y el material
computacional que acompañan a \nolinkurl{PAPER_V_preprint_v1.0_en.md} y
\nolinkurl{PAPER_V_preprint_v1.0_es.md}. El Paper V estudia
empaquetamientos finitos y realizaciones físicas generales. El Paper IV
demuestra la cota para todo orden salvo una constante aditiva absoluta y
el máximo exacto eventual descritos en el artículo principal. La
aplicación IV--V separada no es una premisa de los resultados generales
de V.

\subsection{1. Fuentes canónicas e interfaz
pública}\label{fuentes-canuxf3nicas-e-interfaz-puxfablica}

El punto de entrada canónico es \texttt{PaperV.lean}. Su clausura de
importaciones públicas contiene 122 módulos locales. La inclusión de
archivos, su compilación, su importación pública y la disponibilidad de
los teoremas se comprueban por separado: el build agregado alcanza los
módulos de demostración, los gates de enunciados emiten tipos exactos y
dependencias axiomáticas, y la auditoría completa recorre las
declaraciones de esa misma clausura. El paquete también contiene gates
independientes; son programas de validación, no hipótesis adicionales de
los teoremas públicos.

El entorno es \texttt{leanprover/lean4:v4.35.0}, con Mathlib fijada en
\nolinkurl{56a005855ea17e0b0958f73ec801ae140b0a1679}. El archivo de fuentes
es \nolinkurl{PAPER_V_preprint_v1.0_FORMAL_SOURCE_435.zip}; su
resumen criptográfico exacto y el manifiesto por archivo están
registrados en \nolinkurl{FREEZE_VALIDATION.json}. Las cabeceras reconocen
a Juan Pablo Traverso Gianini por la formalización e integración y
conservan las atribuciones matemáticas y los créditos de fuentes y
adaptación existentes. No reasignan la autoría de los resultados
clásicos.

{\def\LTcaptype{none} % do not increment counter
\begin{longtable}[]{@{}
  >{\raggedright\arraybackslash}p{(\linewidth - 2\tabcolsep) * \real{0.3500}}
  >{\raggedright\arraybackslash}p{(\linewidth - 2\tabcolsep) * \real{0.6500}}@{}}
\toprule\noalign{}
\begin{minipage}[b]{\linewidth}\raggedright
Comprobación
\end{minipage} & \begin{minipage}[b]{\linewidth}\raggedright
Alcance final
\end{minipage} \\
\midrule\noalign{}
\endhead
\bottomrule\noalign{}
\endlastfoot
Raíz pública & \texttt{PaperV} \\
Build agregado & Código 0; 2582 tareas \\
Auditoría completa de declaraciones & Código 0; 2838 declaraciones, 122
módulos, 5547 visitadas \\
Gates de tipos exactos y axiomas & 10 gates; 403 consultas de tipos y
431 consultas de axiomas \\
Axiomas no estándar & Cero \\
Unión de axiomas fundacionales & \texttt{propext},
\texttt{Classical.choice}, \texttt{Quot.sound} \\
Correspondencias principales y de interfaces auxiliares & 37 y 11 \\
\end{longtable}
}

El archivo de fuentes contiene las fuentes locales de Lean comprobadas,
la configuración del paquete, los textos legales y un manifiesto
relativo por archivo. El directorio separado
\nolinkurl{02_validation/PREPRINT_1_0/} contiene los logs finales, los
recibos de salida y la clausura pública exacta. Antes de recorrer las
declaraciones, la auditoría comprueba que todos los módulos enumerados
estén presentes en el entorno importado. Las atribuciones matemáticas y
los créditos de las fuentes se conservan en el desarrollo canónico.

\subsection{2. Arquitectura e
hipótesis}\label{arquitectura-e-hipuxf3tesis}

Los módulos de sistemas de conflicto y copias acotadas definen conjuntos
factibles, defectos de colorabilidad y desigualdades de transferencia.
Los módulos de certificados de recursos conectan relajaciones de
capacidad reales y certificados primal--dual con multiconjuntos de
ocurrencias. Los módulos de realizadores y anfitriones convierten una
clase de color seleccionada en un empaquetamiento físico; el argumento
no toma la unión de clases factibles independientes. Los módulos de
copias, perfiles puros y promoción mixta controlan los recursos físicos
compartidos. Los módulos de frontera unitaria proporcionan restricción,
completación y transporte dual. La sensibilidad finita y la contabilidad
en árboles finitos conservan su propio alcance explícito.

Varias hipótesis deben quedar visibles junto al mapa de declaraciones.
El parámetro de multiplicidad satisface \(q\ge1\). El objetivo ponderado
y el número de clases se fijan antes de seleccionar un certificado. El
LP finito acotado de recursos usa conjuntos de recursos no vacíos donde
se requieren; la relajación general permite conjuntos de recursos
vacíos. Las familias de admisibilidad usadas para tomar un mínimo son no
vacías. Un realizador físico debe cumplir las condiciones reales de
disjunción de recursos. Puentes demostrados relacionan las entradas y
los certificados racionales con los óptimos reales. Los resultados de
copias y copias mixtas conservan sus hipótesis de multiplicidad y
recursos compartidos. Los resultados genéricos de sensibilidad no
identifican una familia WSNF no demostrada; las cotas de árboles finitos
no son cotas globales asintóticas de pérdida.

\subsubsection{2.1. Cobertura de los resultados de
apoyo}\label{cobertura-de-los-resultados-de-apoyo}

{\def\LTcaptype{none} % do not increment counter
\Needspace{18\baselineskip}
\begin{longtable}[]{@{}
  >{\raggedright\arraybackslash}p{(\linewidth - 2\tabcolsep) * \real{0.1500}}
  >{\raggedright\arraybackslash}p{(\linewidth - 2\tabcolsep) * \real{0.8500}}@{}}
\toprule\noalign{}
\begin{minipage}[b]{\linewidth}\raggedright
Elemento
\end{minipage} & \begin{minipage}[b]{\linewidth}\raggedright
Resultado comprobado y límite
\end{minipage} \\
\midrule\noalign{}
\endhead
\bottomrule\noalign{}
\endlastfoot
F1 & Programa lineal finito de recursos efectivo, certificados
racionales primal/dual, realizaciones físicas de clases y cota con dos
errores, incluidos puentes entre óptimos racionales y reales \\
F2 & Portafolio finito de certificados de coste mínimo para un objetivo
y número de clases fijos; uniones físicas compatibles con recursos
disjuntos \\
F3 & Asignaciones voraces de anfitriones por suma de grados, validez
física de recursos, pesos racionales heredados y evitación de un
empaquetamiento de fondo fijo \\
F4 & Promedio ponderado de factores emparejados y completación física de
los supervivientes, incluido el caso de base vacía \\
F5 & Restricción/extensión exacta de frontera unitaria, levantamiento
dual explícito, fórmulas de óptimos enteros y fraccionales y
conservación de la brecha; ganancias originales no negativas, ganancias
centrales reducidas posiblemente con signo \\
F6 & Promoción física mixta de triángulos/K4, conservación de
capacidades, identidades de valor y una secuencia acotada de
normalización que conserva la racionalidad \\
F7 & Sensibilidad genérica de politopos finitos, continuidad, extensión
a familias densas y estabilidad de brechas finitas; no se afirma ninguna
identificación WSNF ni clasificación no demostrada \\
F8 & La contabilidad comprobada de árboles finitos conserva su holgura
fraccional y sus pesos exactos; no es una estimación global de pérdida
asintótica \\
\end{longtable}
}

\subsection{3. SEAL: implementación, garantías y certificados
finitos}\label{seal-implementaciuxf3n-garantuxedas-y-certificados-finitos}

Los resultados matemáticos permanecen en la sección 8 del artículo
principal. La Proposición 8.1 establece la factibilidad de las
operaciones especificadas, la Proposición 8.2 acota las reparaciones de
mejora aceptadas y el Teorema 8.3 proporciona los regímenes de
anfitriones exactos o garantizados allí indicados. Sus declaraciones
formales aparecen en la sección 5 de este documento. Lean verifica esas
interfaces matemáticas bajo sus hipótesis; no verifica la implementación
Python, JavaScript o C++ de extremo a extremo ni su tiempo polinomial.

\subsubsection{3.1. Implementación y modelos
actuales}\label{implementaciuxf3n-y-modelos-actuales}

El software se distribuye en \nolinkurl{03_reproducibility/SEAL/}. Su
directorio \texttt{current/} contiene las rutinas Python BASE, el
programa v3 y los puertos JavaScript/C++ de BASE. \texttt{replay/}
contiene entradas de prueba completas, salidas guardadas, verificadores
exactos y tablas de resultados actuales. El manifiesto fija los bytes
exactos de las fuentes descritas aquí.

Las variantes BASE usan cardinalidad estándar, pesos no negativos de
pares, capacidades de grado o un potencial registrado, según el modelo
de entrada especificado. El programa v3 se restringe a pesos unitarios y
capacidad uno. Se rechazan las llamadas v3 ponderadas o de capacidad
superior no admitidas; las variantes BASE correspondientes siguen siendo
interfaces separadas. Los puertos exponen las familias físicas
seleccionadas además de sus valores objetivos, lo que permite comprobar
independientemente la factibilidad en cada lenguaje ensayado.

La operación N1 actual es transaccional. Libera provisionalmente una
asignación que bloquea, comprueba su relocalización en otro anfitrión
compatible y después comprueba la inserción del par descubierto en el
anfitrión original. Sólo confirma un movimiento factible con ganancia
objetiva estrictamente positiva. Si falla, restaura el estado anterior
completo. Por ello, deben especificarse el objetivo y las restricciones
de capacidad aplicables antes de interpretar una ejecución.

\Needspace{11\baselineskip}
\begin{verbatim}
R := a feasible host assignment
for an uncovered pair e and a compatible host j:
    choose a blocking assignment (f, j)
    save the current state and release (f, j)
    try relocating f to another compatible host
    try inserting e into j
    commit only if feasible and the objective strictly increases
    otherwise restore the saved state
\end{verbatim}

Las rutinas BASE y de respaldo de v3 implementan N1. La rutina separada
\texttt{reference\_v1/} implementa N1 y N2 para el modelo sin pesos y de
capacidad uno; no implementa todas las variantes BASE. En esa
referencia, la selección por escasez va seguida de una búsqueda
exhaustiva en los entornos finitos, que se reinicia tras aceptar una
mejora. La auditoría LP opcional se ejecuta después de la construcción.

La rama bipartita de v3 valida la asignación propia desde listas que
obtiene antes de utilizarla. No se afirma que el ejecutable construya
una asignación en todo el rango del teorema BKW. Una asignación
individual satisfactoria o la coincidencia de cotas inferior y superior
proporciona su propio certificado finito. Los regímenes matemáticos
especiales exactos no implican una garantía general de aproximación para
la búsqueda local.

\Needspace{14\baselineskip}
\subsubsection{3.2. Resultados finitos
actuales}\label{resultados-finitos-actuales}

{\def\LTcaptype{none} % do not increment counter
\begin{longtable}[]{@{}
  >{\raggedright\arraybackslash}p{(\linewidth - 2\tabcolsep) * \real{0.5000}}
  >{\raggedright\arraybackslash}p{(\linewidth - 2\tabcolsep) * \real{0.5000}}@{}}
\toprule\noalign{}
\begin{minipage}[b]{\linewidth}\raggedright
Comprobación del conjunto actual
\end{minipage} & \begin{minipage}[b]{\linewidth}\raggedright
Resultado verificado
\end{minipage} \\
\midrule\noalign{}
\endhead
\bottomrule\noalign{}
\endlastfoot
Empaquetamientos físicos Python & 110: 72 BASE y 38 v3 \\
Empaquetamientos físicos JavaScript y C++ & 54 en cada lenguaje; todos
los valores objetivos coinciden con Python \\
Duales racionales completos de cardinalidad & 38; 149,819 desigualdades
de objetos y 47,015 filas comprobadas \\
Óptimos de cardinalidad certificados por cotas iguales & 12: seis casos
de distribución, dos testigos y cuatro casos mayores \\
Intervalos certificados positivos & 26: doce gaps de distribución de
1--3 y catorce gaps de casos mayores de 1--40 \\
Controles negativos independientes & Siete rechazados \\
\end{longtable}
}

Las pruebas BASE usan dieciocho entradas guardadas en cada uno de cuatro
modos. Las entradas ponderadas usan pesos en \texttt{\{1,2,5,10\}}; las
de mayor capacidad asignan capacidad dos a cada anfitrión. Cada puerto
comprueba los modos estándar, ponderado y de capacidad dos sobre esas
entradas. Las pruebas v3 usan pesos unitarios y capacidad uno: dieciocho
instancias de distribución, dos testigos bipartitos y dieciocho
configuraciones grandes guardadas. Las comprobaciones de factibilidad y
objetivo ponderadas y de capacidad dos no proporcionan certificados de
optimalidad para esos modelos. Los valores actuales completos por
instancia se registran en \nolinkurl{replay/CURRENT_RESULTS.csv}.

Un testigo físico N1 explícito pasa de un par seleccionado a dos. La
regresión transaccional finita adicional comprueba 685 estados iniciales
factibles y 27,708 intentos de operación, con 3,180 inserciones
aceptadas y 24,528 intentos rechazados que restauran el estado. Esas
cifras corresponden a dos sistemas de anfitriones especificados,
capacidades uno/dos y ganancias cero/uno; no son afirmaciones
universales de corrección.

El verificador de reproducción reconstruye la pertenencia a los
anfitriones, la unicidad global de pares, las capacidades de grado por
anfitrión, las cardinalidades y los objetivos ponderados desde las
especificaciones completas de entrada. Las pruebas transaccionales
comprueban mejoras N1 aceptadas, la restauración tras intentos fallidos,
el comportamiento con peso cero y los modelos de capacidad indicados.
Son comprobaciones finitas ejecutadas, no demostraciones universales de
corrección del programa.

Para cada dual racional guardado, reconstruye todos los objetos y filas
de restricciones, comprueba los precios no negativos y la cobertura, y
recalcula exactamente el objetivo. Un empaquetamiento factible da una
cota inferior; la parte entera de una cota superior racional válida da
una cota superior entera. La coincidencia de cotas certifica el óptimo
de esa salida individual, mientras que un intervalo positivo registra
incertidumbre. No se usa como certificado un mero estado del solver de
punto flotante ni un objetivo numérico.

La comparación entre los tres lenguajes se limita a los modelos de
entrada y las salidas guardadas. No demuestra equivalencia general, y
los tiempos no son premisas matemáticas. El informe de resultados
distribuido con el software conserva los certificados exactos actuales,
los hashes de fuentes y los límites de alcance.

El conjunto adicional de referencia de capacidad uno contiene treinta
instancias guardadas y 147,111 objetos anfitrión--par. Catorce
intervalos certifican óptimos enteros; dieciséis conservan incertidumbre
positiva, de ancho entero a lo sumo 37. Seis óptimos pequeños tienen
enumeración independiente. Se rechazan cuatro certificados alterados y
dos pruebas explícitas ejercitan N1 y N2. Doce casos compactos habilitan
reparaciones; dieciocho casos mayores sólo ejecutan selección y
auditoría dual. Las restricciones de conjuntos impares se generan para
paletas de tamaño a lo sumo nueve; omitir conjuntos impares mayores
puede debilitar una cota sin invalidar un dual comprobado. Este conjunto
no implica exactitud universal de la relajación.

\subsubsection{3.3. Comprobaciones de certificados
escalares}\label{comprobaciones-de-certificados-escalares}

Certo 0.34.0 comprueba cinco consecuencias escalares codificadas. Hacen
explícitas las hipótesis aritméticas y proporcionan comprobaciones
reproducibles; no demuestran la existencia de un grafo, un realizador
físico ni un empaquetamiento factible.

{\def\LTcaptype{none} % do not increment counter
\begin{longtable}[]{@{}
  >{\raggedright\arraybackslash}p{(\linewidth - 4\tabcolsep) * \real{0.3333}}
  >{\raggedright\arraybackslash}p{(\linewidth - 4\tabcolsep) * \real{0.3333}}
  >{\raggedright\arraybackslash}p{(\linewidth - 4\tabcolsep) * \real{0.3333}}@{}}
\toprule\noalign{}
\begin{minipage}[b]{\linewidth}\raggedright
Obligación
\end{minipage} & \begin{minipage}[b]{\linewidth}\raggedright
Papel matemático
\end{minipage} & \begin{minipage}[b]{\linewidth}\raggedright
Condición de reproducción
\end{minipage} \\
\midrule\noalign{}
\endhead
\bottomrule\noalign{}
\endlastfoot
\nolinkurl{two_error_transfer} & Combinar pérdidas de relajación y
realización & Requiere SMT \\
\nolinkurl{integer_dual_stop} & Parar cuando coinciden las cotas inferior
entera y superior dual & Requiere SMT \\
\nolinkurl{composed_realization_errors} & Sumar los presupuestos de error
de realización indicados & Requiere SMT \\
\nolinkurl{improving_repairs_budget} & Acotar las reparaciones de mejora
aceptadas & Certificado sin solver \\
\nolinkurl{gap_residual_and_resource_defect} & Separar residuo y
defecto de recurso & Requiere SMT \\
\end{longtable}
}

Un sexto certificado, independiente del solver, es un control negativo
de una implicación escalar deliberadamente insuficiente:
\(\nu=\mathrm{cov}=0\) y un valor de relajación \(1/2\) satisfacen esas
premisas codificadas y dejan un gap de relajación positivo. No es un
contraejemplo de grafos ni se afirma que proceda de un sistema de
anfitriones. En particular, no contradice las cotas superiores de
recursos conservadas por el modelo real de SEAL. Cada comprobación
vincula el hash de la especificación y la obligación; se rechaza la
comprobación contra una obligación diferente. \texttt{RESULTS.json}
registra qué comprobaciones requieren un solver SMT.

\subsection{4. Reproducción}\label{reproducciuxf3n}

Extraer \nolinkurl{PAPER_V_preprint_v1.0_FORMAL_SOURCE_435.zip} y
ejecutar desde la raíz del paquete con las dependencias fijadas:

\Needspace{9\baselineskip}
\begin{verbatim}
lake build PaperV
lake env lean PaperVManuscriptGate.lean
lake env lean PaperVSupplementGate.lean
lake env lean PaperVImprovementsGate.lean
lake env lean PaperVReleaseGate.lean
lake env lean FreezeFullAudit.lean
\end{verbatim}

\nolinkurl{CHECKED_PUBLICATION_GATES_435.json} registra los diez gates,
los comandos exactos, las declaraciones consultadas y sus dependencias
axiomáticas. \nolinkurl{MANUSCRIPT_MAP.json} y
\nolinkurl{SUPPLEMENT_MAP.json} dan las correspondencias exactas por
enunciado. \nolinkurl{SOURCE_CLOSURE.json} identifica los módulos públicos
alcanzados por el agregado. La auditoría no admite nuevos axiomas
matemáticos.

Desde \nolinkurl{03_reproducibility/SEAL/}, comprobar las salidas
guardadas con la biblioteca estándar de Python:

\Needspace{4\baselineskip}
\begin{verbatim}
python replay/scripts/check_outputs.py --source-root current
\end{verbatim}

Para regenerar el conjunto actual, usar las dependencias y entornos de
ejecución fijados en su README y ejecutar:

\Needspace{6\baselineskip}
\begin{verbatim}
python replay/scripts/run_suite.py --source-root current
python replay/scripts/run_ports.py --source-root current
python replay/scripts/check_outputs.py --source-root current
\end{verbatim}

Los testigos adicionales de referencia y los certificados escalares
tienen comandos separados:

\Needspace{5\baselineskip}
\begin{verbatim}
python reference_v1/check_certificate.py
python certo_v/run_certo.py
\end{verbatim}

Ejecutar Python normalmente, sin \texttt{-O}, para mantener las
aserciones. Comprobar los testigos guardados no requiere el solver
numérico que generó sus precios. La generación y las comprobaciones de
Certo que necesitan SMT requieren sus dependencias documentadas; generar
los puertos requiere Node y un compilador C++17, que puede especificarse
con \texttt{-\/-cxx}. El manifiesto del software vincula fuentes,
entradas, vectores completos, salidas y registros de validación.

\subsection{5. Mapa de declaraciones por
enunciado}\label{mapa-de-declaraciones-por-enunciado}

Las 37 correspondencias siguientes identifican las declaraciones
públicas. Los tipos exactos, las hipótesis y las conclusiones separadas
están en las fuentes y en la salida de los gates; este índice no suprime
las precisiones de la sección 2.

{\def\LTcaptype{none} % do not increment counter
\begin{longtable}[]{@{}
  >{\raggedright\arraybackslash}p{(\linewidth - 2\tabcolsep) * \real{0.2800}}
  >{\raggedright\arraybackslash}p{(\linewidth - 2\tabcolsep) * \real{0.7200}}@{}}
\toprule\noalign{}
\begin{minipage}[b]{\linewidth}\raggedright
Elemento del artículo principal
\end{minipage} & \begin{minipage}[b]{\linewidth}\raggedright
Declaraciones canónicas
\end{minipage} \\
\midrule\noalign{}
\endhead
\bottomrule\noalign{}
\endlastfoot
Definición 2.1 & \nolinkurl{PaperV.ConflictSystem};
\texttt{PaperV.CSFeasible}; \texttt{PaperV.nu} \\
Definición 2.2 & \texttt{PaperV.Copies}; \texttt{PaperV.Hgraph};
\nolinkurl{PaperV.IsQMultiset} \\
Definición 2.3 & \texttt{PaperV.Dq}; \texttt{PaperV.oct};
\nolinkurl{PaperV.Dq_one_eq} \\
Definición 2.4 & \nolinkurl{PaperV.Manuscript.ResourceRelaxation};
\nolinkurl{PaperV.ResourceCertificate.System};
\nolinkurl{PaperV.ResourceCertificate.System.incidence} \\
Definición 2.5 & \nolinkurl{PaperV.Manuscript.ResourceRelaxation};
\nolinkurl{PaperV.Manuscript.ResourceRelaxation.Optimum} \\
Lema 3.1 & \texttt{PaperV.splitting} \\
Corolario 3.2 & \nolinkurl{PaperV.nu_ge_ceil} \\
Corolario 3.2a & \nolinkurl{PaperV.cover_witness_defect};
\nolinkurl{PaperV.triangle_cover_defect} \\
Corolario 3.2b &
\nolinkurl{PaperV.Manuscript.minimum_coloring_rigidity_iff} \\
Teorema 3.3 &
\nolinkurl{PaperV.Manuscript.resource_half_integral_transfer} \\
Teorema 3.4 & \nolinkurl{PaperV.Manuscript.resource_integral_transfer} \\
Teorema 3.4w & \nolinkurl{PaperV.transfer_q_weighted} \\
Lema 3.5 & \nolinkurl{PaperV.parity_localization} \\
Definición 4.1 & \nolinkurl{PaperV.ClassRealizer};
\texttt{PaperV.RealCert}; \nolinkurl{PaperV.RealCert.cost} \\
Teorema 4.2 & \nolinkurl{PaperV.realizable_class_transfer};
\nolinkurl{PaperV.realizable_transfer_cert} \\
Definición 4.3 & \texttt{PaperV.wRDq} \\
Teorema 4.4 & \nolinkurl{PaperV.wRDq_identity_eq_wDq};
\nolinkurl{PaperV.wRDq_le_wDq} \\
Ejemplo 4.5 & \nolinkurl{PaperV.StarExample.intermediate_defect} \\
Corolario 4.6 &
\nolinkurl{PaperV.ResourceCertificate.two_error_real_optimum};
\nolinkurl{PaperV.ResourceCertificate.resource_gap_identity};
\nolinkurl{PaperV.ResourceCertificate.exists_real_optimum} \\
Proposición 5.1 &
\nolinkurl{PaperV.Host.csFeasible_hostCS_of_properOn} \\
Teorema 5.2 & \nolinkurl{PaperV.Galvin.galvin_multigraph_max_degree};
\nolinkurl{PaperV.Galvin.galvin_bkw_local} \\
Corolario 5.3 & \nolinkurl{PaperV.Host.galvin_host_completion_exact};
\nolinkurl{PaperV.Host.bkw_host_completion_exact};
\nolinkurl{PaperV.Host.galvin_lifted_transfer};
\nolinkurl{PaperV.Host.galvin_local_lifted_transfer} \\
Ejemplo 5.5 & \nolinkurl{PaperV.Strictness.strict_separation};
\nolinkurl{PaperV.Strictness.nu_hostCS_kr} \\
Afirmación auxiliar 6.1a &
\nolinkurl{PaperV.copy_integral_edgeLoad_le_two} \\
Teorema 6.1 & \nolinkurl{PaperV.copy_fractional} \\
Teorema 6.2 & \nolinkurl{PaperV.copy_integral_for_optimal_packings};
\nolinkurl{PaperV.copy_integral_for_packings};
\nolinkurl{PaperV.triangleImageMultiset} \\
Lema 6.2k & \nolinkurl{PaperV.third_subset_union};
\nolinkurl{PaperV.k4_of_three_triangles};
\nolinkurl{PaperV.triCS_clique_card_le_four} \\
Corolario 6.2a & \nolinkurl{PaperV.copy_triple_bound};
\texttt{PaperV.tight\_oct} \\
Proposición 7.1 & \nolinkurl{PaperV.PureProfile.prop43_nu};
\nolinkurl{PaperV.PureProfile.prop43_muMat} \\
Proposición 7.1a & \nolinkurl{PaperV.PureProfile.nu_ge_frac_muMatDeg};
\nolinkurl{PaperV.PureProfile.nu_ge_two_thirds_muMatDeg};
\nolinkurl{PaperV.PureProfile.nu_ge_muMatDeg_two} \\
Proposición 7.1b & \nolinkurl{PaperV.PureProfile.muMatDeg_eq_min};
\nolinkurl{PaperV.PureProfile.nu_eq_muMatDeg_of_even};
\nolinkurl{PaperV.PureProfile.nu_eq_frac_muMatDeg_of_odd} \\
Teorema 7.3 & \nolinkurl{PaperV.Mixed.mixed_copy_fractional};
\nolinkurl{PaperV.mixed_copy_integral} \\
Corolario 7.3a &
\nolinkurl{PaperV.Mixed.mixed_copy_integral_of_roundingGap} \\
Proposición 8.1 & \nolinkurl{PaperV.feasible_run};
\nolinkurl{PaperV.SealBridge.sealFeasible_of_properOn} \\
Proposición 8.2 & \nolinkurl{PaperV.accepted_repairs_le_cov};
\nolinkurl{PaperV.accepted_repairs_le_audit} \\
Teorema 8.3 & \nolinkurl{PaperV.SealBridge.seal_galvin_regime};
\nolinkurl{PaperV.SealBridge.seal_bkw_regime};
\nolinkurl{PaperV.SealBridge.seal_konig_guarantee} \\
Proposición 9.3 &
\nolinkurl{PaperV.Steiner.exists_quadratic_defect_frac} \\
\end{longtable}
}

La interfaz del Teorema 6.2 recibe los dos empaquetamientos copiados
factibles especificados y construye su multiconjunto imagen real; el
adaptador para empaquetamientos óptimos conserva esas mismas familias.
La interfaz del Teorema 7.3 también utiliza los datos reales de la copia
mixta. Se conservan adaptadores existenciales por conveniencia, pero no
sustituyen estas correspondencias.

\subsection{6. Interfaces matemáticas
auxiliares}\label{interfaces-matemuxe1ticas-auxiliares}

Las once correspondencias siguientes localizan los resultados auxiliares
en el proyecto de fuentes entregado. Sus etiquetas identifican el
inventario matemático de las fuentes; los tipos exactos determinan todas
las hipótesis. No hace falta un suplemento separado no distribuido para
localizar estas interfaces.

{\def\LTcaptype{none} % do not increment counter
\begin{longtable}[]{@{}
  >{\raggedright\arraybackslash}p{(\linewidth - 2\tabcolsep) * \real{0.2800}}
  >{\raggedright\arraybackslash}p{(\linewidth - 2\tabcolsep) * \real{0.7200}}@{}}
\toprule\noalign{}
\begin{minipage}[b]{\linewidth}\raggedright
Elemento conservado
\end{minipage} & \begin{minipage}[b]{\linewidth}\raggedright
Declaraciones canónicas
\end{minipage} \\
\midrule\noalign{}
\endhead
\bottomrule\noalign{}
\endlastfoot
Proposición 3.6 & \nolinkurl{PaperV.Dq_antitone};
\nolinkurl{PaperV.Dq_le_oct}; \nolinkurl{PaperV.Dq_mono_left} \\
Proposición 3.7 & \nolinkurl{PaperV.Dq_eq_zero_iff};
\nolinkurl{PaperV.oct_eq_zero_iff};
\nolinkurl{PaperV.Dq_eq_zero_iff_chromatic};
\nolinkurl{PaperV.Dq_pos_iff_lt_chromatic} \\
Proposición 3.8 & \nolinkurl{PaperV.Dq_ge_card_sub};
\nolinkurl{PaperV.defect_tight_all_max};
\nolinkurl{PaperV.defect_eq_of_survivors};
\nolinkurl{PaperV.Manuscript.minimum_coloring_rigidity_iff} \\
Proposición 3.9 & \texttt{PaperV.Dq\_sum}; \texttt{PaperV.oct\_sum} \\
Proposición 3.10 & \nolinkurl{PaperV.Dq_lineGraph_eq_zero_iff};
\nolinkurl{PaperV.Dq_lineGraph_pos_iff};
\nolinkurl{PaperV.Dq_one_lineGraph};
\nolinkurl{PaperV.profile_zero_set_not_values} \\
Corolario 3.11 & \nolinkurl{PaperV.TerminalPair.p4_133_existence_sharp};
\nolinkurl{PaperV.TerminalPair.p4_133_existence_coarse} \\
Teorema 7.4 &
\nolinkurl{PaperV.PairedFactors.exists_paired_completion} \\
Proposición 7.5 & \nolinkurl{PaperV.FiniteLP.gap_identity};
\nolinkurl{PaperV.FiniteLP.weak_duality};
\nolinkurl{PaperV.FiniteLP.complementarity_of_equal_value};
\nolinkurl{PaperV.RationalFarkas.farkas};
\nolinkurl{PaperV.RationalFarkas.not_sat_of_cert};
\nolinkurl{PaperV.FiniteLP.exists_equal_primal_dual} \\
Corolario 5.4 & \nolinkurl{PaperV.Host.weighted_host_selection};
\nolinkurl{PaperV.Host.weighted_host_selection_nu};
\nolinkurl{PaperV.Host.konig_weighted_host_selection};
\nolinkurl{PaperV.Host.galvin_heavy_single_host} \\
Lema C.1 &
\nolinkurl{PaperV.UnitBoundary.completion_optimal_unconditional};
\nolinkurl{PaperV.UnitBoundary.unitBoundary_completion_optimal};
\nolinkurl{PaperV.UnitBoundary.exists_optimal_normPack} \\
Lema C.2 & \nolinkurl{PaperV.UnitBoundary.exists_endpoint_repacking};
\nolinkurl{PaperV.UnitBoundary.endpoint_repacking_card_bound};
\nolinkurl{PaperV.UnitBoundary.endpoint_repacking_nonloss} \\
\end{longtable}
}

El Lema C.1 tiene un adaptador incondicional de la fórmula del óptimo
sin hipótesis de universo habitado ni de representantes; su versión con
testigo normalizado conserva representantes explícitos. El Lema C.2
proporciona una familia real de representantes y un empaquetamiento
físico, incluido el universo vacío; el enunciado de transporte conserva
sus datos de empaquetamiento normalizado.

\subsection{7. Alcance del registro
final}\label{alcance-del-registro-final}

El manuscrito principal contiene un párrafo de formalización; este
companion reúne la arquitectura de fuentes, las correspondencias
exactas, el informe computacional y los comandos de reproducción. La
comprobación formal, la comprobación de certificados finitos, la
prioridad bibliográfica y la revisión humana externa son obligaciones
diferentes. La aplicación IV--V separada tiene su propio proyecto y
registro de validación. El Paper V no importa la cota global de esa
aplicación ni depende del Paper VI. Este companion acompaña el Preprint
1.0. La comprobación por el kernel y la verificación de certificados
finitos no constituyen revisión humana externa ni certificación de
prioridad histórica.

\Needspace{18\baselineskip}
\subsection{8. Reproducibilidad de los contraejemplos
finitos}\label{reproducibilidad-de-los-contraejemplos-finitos}

El directorio \nolinkurl{03_reproducibility/counterexamples/} proporciona
\nolinkurl{bk_completion_certificates.py} y el certificado BK completo
guardado. Desde ese directorio se ejecuta
\texttt{python\ bk\_completion\_certificates.py\ -\/-verify-only}, con
las aserciones activadas. El comprobador usa sólo la biblioteca
estándar, valida el dual racional de valor 19/2, enumera los 70
empaquetamientos óptimos de nueve elementos y comprueba el parámetro
residual beta=2. Se rechaza un control negativo que altera un precio
dual. La búsqueda de certificados sin \texttt{-\/-verify-only} utiliza
SciPy de manera opcional; verificar el certificado guardado no lo
requiere. El número exacto 70 es un resultado computacional, no una
declaración demostrada por el kernel de Lean. Declaraciones comprobadas
por el kernel establecen por separado el óptimo BK y las cotas
pertinentes, así como los hechos enunciados sobre CLON-2, STAR-T2 y
Steiner.

\subsection{9. Licencias y atribución}\label{licencias-y-atribuciuxf3n}

Copyright (c) 2026 Juan Pablo Traverso Gianini. Los materiales escritos
se publican bajo CC BY-NC 4.0. Las contribuciones del autor a Lean y al
software se publican bajo Apache 2.0; las dependencias y contribuciones
de terceros conservan sus propias licencias y avisos. Las cabeceras de
las fuentes y los textos legales adjuntos explicitan esta distinción. Se
conservan las atribuciones matemáticas clásicas; reconocer una
formalización no reatribuye el teorema subyacente.

\end{document}
